\documentclass[aspectratio=169,usenames,dvipsnames]{beamer}
\beamertemplatenavigationsymbolsempty
% \documentclass{beamer}

% UNCOMMENT FOR HANDOUTS
% \usepackage{pgfpages}
% \pgfpagesuselayout{2 on 1}[landscape,letterpaper,border shrink=5mm]
% \setbeameroption{show notes}

\mode<presentation>
\usetheme{Boadilla}
\usecolortheme{spruce}
\usepackage{../../LizStyle,ulem}
\usepackage{color}
\usepackage[color,matrix,arrow]{xy}
% \input xy
% \xyoption{all}
\usepackage{tikz}
\usepackage{tikz-cd}
\tikzset{commutative diagrams/diagrams={ampersand replacement=\&}}

\AtBeginSection{\frame{\sectionpage}}


% ==========Command for putting comments and removing them=======
% Visible:
\newcommand{\InstructorNotes}[1]{\textit{\color{orange}#1}}
\newcommand{\InstructorList}[1]{
\begin{itemize}
	\color{orange}
	#1
\end{itemize}
}
% Removed:
% \renewcommand{\InstructorNotes}[1]{}
% \renewcommand{\InstructorList}[1]{}


\newcommand{\bvfill}{\vskip0pt plus 1filll}

\usepackage{multimedia}
\usepackage{graphbox}

% \usepackage{algpseudocode}


\newcommand{\Reeb}{\textbf{Reeb}}
\newcommand{\Rtopc}{\textbf{Rtopc}}
% \newcommand{\RR}{\mathcal{R}}

\newcommand{\Set}{\textbf{Set}}
\newcommand{\Vect}{\textbf{Vect}}
\newcommand{\Top}{\textbf{Top}}
\newcommand{\Int}{\textbf{Int}}


\graphicspath{ {Figures/} {../Figures/}{../../Figures/} }



\begin{document}
%\pdfoptionpdfminorversion 5
%\logo{\includegraphics[height=1.5cm]{dukeshield.pdf}}
\title[]{Ch 10.1 - 10.3: Neural Nets}
\subtitle{Lecture 25 - CMSE 381}
\date{Mon, Apr 18, 2022}



\institute[MSU-CMSE]{Michigan State University\\
			::\\
          Dept of Computational Mathematics, Science \& Engineering}
\author[Dr.~Munch]{Prof.~Elizabeth Munch}

\begin{frame}
 \titlepage


\end{frame}
%--------------------------
\begin{frame}[c]{Announcements }

\begin{columns}
\column{.35\textwidth}
\centering

	\textbf{Last time:}
	\begin{itemize}
	\item SVM Lab
	\end{itemize}

	\vspace{.2in}

	\textbf{This lecture:}
	\begin{itemize}
		\item Feed Forward Neural Nets
		\item CNNs
\end{itemize}


\column{.55\textwidth}
\centering

\textbf{Announcements:}
	\begin{itemize}
	\item Mon 4/18: 10.1-3
	\item Weds 4/20: 10.4-5
	\item Fri 4/22: Lab: 10.9
	\begin{itemize}
		\item HW \#10 Due Friday
	\end{itemize}
	\item Week of 4/25: Ch 12 Clustering
	\item Final: May 4, 5:45 - 7:45 PM, Erickson 226

	\end{itemize}
\end{columns}

\end{frame}
%--- Next Frame ---%


%==========
\section{Neural Nets}
%==========

\begin{frame}[t]{The idea}

\begin{columns}
\column{.45\textwidth}
\centering
\InstructorList{
  \item Draw input of vector $p$ variables $X = (X_1,\cdots,X_p)$
	\item Draw black box nonlinear function $f(X)$
	\item Spit out response prediction $Y$
	\\
	\item We've had nonlinear stuff before like trees, boosting, GANs
	\item Difference here is the structure
}

\column{.45\textwidth}
\centering

\end{columns}

\end{frame}
%--- Next Frame ---%
\begin{frame}[t]{Feed Forward Neural Network: The cartoon}

\begin{columns}
\column{.45\textwidth}
\centering
\includegraphics[width = \textwidth, align = c]{Fig10-1}
\column{.45\textwidth}
\centering
\InstructorList{
  \item Input layer is starting data
	\item Each arrow means taking combo of those values with weights to get next value
	\item Here there is one hidden layer then the output
	\item Get to pick how many hidden units $K$, here we have $K=5$
}

\end{columns}

\end{frame}
%--- Next Frame ---%
\begin{frame}[t]{A very simple example}
	{Computing $A_k$ for $(1,0)$}

\includegraphics[width = .4\textwidth, align = c]{Eq10-2}

\begin{columns}
\column{.45\textwidth}
\centering
\includegraphics[width = \textwidth, align = c]{DL-simple-ex}
\column{.45\textwidth}
\centering
\InstructorNotes{
\includegraphics[width = .5\textwidth, align = c]{DL-simple-ex-Input10}
\includegraphics[width = .5\textwidth, align = c]{DL-simple-ex-Input10-Soln}
}
\InstructorList{
  \item $A_k$ called activations
	\item Pick an activation function $g(z)$
}

\end{columns}

\end{frame}
%--- Next Frame ---%


\begin{frame}[t]{A very simple example}
	{Computing $Y$ for $(1,0)$}


$\displaystyle{f(X) = \beta_0 + \sum_{k=1}^K \beta_k A_k}$

\begin{columns}
\column{.45\textwidth}
\centering
\includegraphics[width = \textwidth, align = c]{DL-simple-ex}
\column{.45\textwidth}
\centering
\InstructorNotes{
\includegraphics[width = .5\textwidth, align = c]{DL-simple-ex-Input10}
\includegraphics[width = .5\textwidth, align = c]{DL-simple-ex-Input10-Soln}
}
\end{columns}

\end{frame}
%--- Next Frame ---%


\begin{frame}[t]{A very simple example}
	{Computing $Y$ for $(0,1)$}


\includegraphics[width = .4\textwidth, align = c]{Eq10-2}
\qquad
$\displaystyle{f(X) = \beta_0 + \sum_{k=1}^K \beta_k A_k}$

\begin{columns}
\column{.45\textwidth}
\centering
\includegraphics[width = \textwidth, align = c]{DL-simple-ex}
\column{.45\textwidth}
\centering
\InstructorNotes{
\includegraphics[width = .5\textwidth, align = c]{DL-simple-ex-Input01}
\includegraphics[width = .5\textwidth, align = c]{DL-simple-ex-Input01-Soln}
}
\end{columns}

\end{frame}
%--- Next Frame ---%

\begin{frame}[t]{Choices for activation function }

\begin{columns}
\column{.45\textwidth}
\centering
Sigmoid:
\begin{equation*}
	g(z) = \frac{e^z}{1+e^z} = \frac{1}{1+e^{-z}}
\end{equation*}
\InstructorList{
  \item Converts linear function in probabilities in [0,1]
}

\column{.45\textwidth}
\centering
ReLU: Rectified linear unit
\begin{equation*}
	g(z) = (z)_+ =
	\begin{cases}
		0 & \text{if }z<0 \\
		z & \text{else.}
	\end{cases}
\end{equation*}
\end{columns}

\end{frame}
%--- Next Frame ---%


\begin{frame}[t]{Now what?}
Choose parameters by minimizing RSS, $\sum_{i=1}^n(y_i - f(x_i))^2$
\begin{columns}
\column{.45\textwidth}
\centering
\textbf{Chosen in advance:}
\InstructorList{
  \item Number of layers (more on that in a moment)
	\item Number of hidden units
	\item Activation function $g(z)$
}


\column{.45\textwidth}
\centering
\textbf{Tuned by the model:}
\InstructorList{
  \item $\omega$
	\item $\beta$
}

\end{columns}

\end{frame}
%--- Next Frame ---%

%==========
\section{Multilayer Neural Networks}
%==========



\begin{frame}[t]{MNIST}

\begin{columns}
\column{.45\textwidth}
\centering
\includegraphics[width = \textwidth, align = c]{Fig10-3}
\column{.45\textwidth}
\centering
\InstructorList{
  \item Goal: Build a model to classify images into their correct digit class
	\item Each image has $p=28 \cdot 28 = 784$ pixels
	\item Each pixel is grayscale value in [0,255]
	\item Data converted into column order
	\item Output represented by one-hot vector $Y = (Y_0,Y_1,\cdots,Y_9)$
	\item 60K training images, 10K test images
}

\end{columns}

\end{frame}
%--- Next Frame ---%
\begin{frame}[t]{Neural network architecture for MNIST}

\begin{columns}
\column{.45\textwidth}
\centering
\includegraphics[width = \textwidth, align = c]{Fig10-4}
\column{.45\textwidth}
\centering
\includegraphics[width = \textwidth, align = c]{Eq10-10}
\includegraphics[width = \textwidth, align = c]{Eq10-11}
\InstructorList{
  \item Two hidden layers. $L_1$ has 256 units; $L_2$ has 128
	\item 10 output variables due to class labeling
}

\end{columns}

\end{frame}
%--- Next Frame ---%

\begin{frame}[t]{More on that architecture}

\begin{columns}
\column{.45\textwidth}
\centering
\includegraphics[width = \textwidth, align = c]{Fig10-4}
\column{.45\textwidth}
\centering
\includegraphics[width = \textwidth, align = c]{Eq10-10}
\includegraphics[width = \textwidth, align = c]{Eq10-11}
\InstructorList{
  \item Superscript denotes layer
	\item $W_1$ denotes entire matrix of weight
	\item In this setting, size is $785 \times 256 = 200,960$ values
	\item 785 instead of 784 to involve intercept term (called bias in this literature)
}

\end{columns}

\end{frame}
%--- Next Frame ---%
\begin{frame}[t]{The last column}

\begin{columns}
\column{.45\textwidth}
\centering

\includegraphics[width = \textwidth, align = c]{Fig10-4}
\column{.45\textwidth}
\centering
\includegraphics[width = \textwidth, align = c]{Eq10-12}
\includegraphics[width = \textwidth, align = c]{Eq10-13}
\InstructorList{
  \item Doing linear combo of inputs to get 10 different outputs
	\item Want this to act like a probability.
	\item Answer is to use the softmax activation function $f_m(X)$
	\item Values are non-negative and sum to 1
	\item Return class with highest probability
}


\end{columns}

\end{frame}
%--- Next Frame ---%
\begin{frame}[t]{Results}

\begin{columns}
\column{.45\textwidth}
\centering
\includegraphics[width = .5\textwidth, align = c]{Fig10-3}
\includegraphics[width = \textwidth, align = c]{Table10-1}
\column{.45\textwidth}
\centering
\InstructorList{
  \item Bottom two are previously covered
	\item Dramatic increase with neural nets
	\item Dropout regularization covered later..... maybe we won't get there
	\item Number of coefficients to train in $W_1$, $W_2$, and $B$ is 235,146; more than 33 times number needed for multinomial logistic regression
	\item $p \approx 4 n$, so need to avoid overfitting
	\item Regularization is important, but we won't get to that.

}

\end{columns}

\end{frame}
%--- Next Frame ---%

%==========
\section{Convolutional Neural Network }
%==========

\begin{frame}[t]{CIFAR100 Data}

	\centering
\includegraphics[width = .6\textwidth, align = c]{Fig10-5}
\begin{columns}
\column{.45\textwidth}
\centering
\InstructorList{
  \item 60,000 images: 50K training, 10K  test
	\item Labels with 20 super classes (e.g.~aquatic mammals)
	\item 5 classes per super class (beaver, dolphin, otter, seal, whale)
	\item Images are 32x32
}

\column{.45\textwidth}
\centering
\footnotesize


\end{columns}

\end{frame}
%--- Next Frame ---%
\begin{frame}[t]{Image channel data}

\begin{columns}
\column{.45\textwidth}
\centering

\InstructorList{
	\item 3 eight bit numbers per pixel (red, green, blue)
	\item Numbers are organized in a \textit{feature map} , which is 3 dimensional array
	\item First two directions are spatial (32x32)
	\item Third is the \textit{channel} axis representing teh three colors
}
\column{.45\textwidth}
\centering
\includegraphics[width = .5\textwidth]{Fig10-8-onlychannels.png}
\end{columns}

\end{frame}
%--- Next Frame ---%
\begin{frame}[t]{CNNs}

\begin{columns}
\column{.45\textwidth}
\centering
\InstructorList{
  \item Idea is to build a neural net that works on this data
	\item Identify low level features on the input image such as small edges, patches of color, etc
	\item Low level features combined into higher-level features (parts of ears, eyes, etc)
	\item Does this using special types of hidden layers: \textit{convolution} layers and \textit{pooling layers}
}

\column{.45\textwidth}
\centering
\includegraphics[width = \textwidth, align = c]{Fig10-6}

\end{columns}

\end{frame}
%--- Next Frame ---%

\begin{frame}[t]{Convolution layer }
	{Convolution Filter}

\begin{columns}
\column{.45\textwidth}
\centering
Original Image:
\begin{equation*}
	\begin{bmatrix}
		a & b & c\\
		d & e & f \\
		g & h & i\\
		j & k & l
	\end{bmatrix}
\end{equation*}

Convolution filter:
\begin{equation*}
	\begin{bmatrix}
		\alpha & \beta \\
		\gamma & \delta
	\end{bmatrix}
\end{equation*}
\column{.45\textwidth}
\centering
Convolved Image
\includegraphics[width = \textwidth, align = c]{CNN-ConvolvedImage.png}
\end{columns}

\end{frame}
%--- Next Frame ---%
\begin{frame}[t]{Convolution Filter Example}

\begin{columns}
\column{.25\textwidth}
\centering
\includegraphics[width = \textwidth, align = c]{CNN-examplematrix}
% \includegraphics[width = .5\textwidth, align = c]{CNN-examplematrix-Mask}

Filter:
\begin{equation*}
	\begin{bmatrix}
		0 & 1 \\
		0 & 1
	\end{bmatrix}
\end{equation*}
\column{.45\textwidth}
\centering

\InstructorList{
\item $\begin{bmatrix}
	1 & 3 \\
	0 & 3
\end{bmatrix}$
\item Convolved image highlights regions that resemble the convolution filter.
\item Note that the result is smaller than original
}
\end{columns}
\end{frame}
%--- Next Frame ---%
\begin{frame}[t]{Convolution filter: Bigger example}

	\centering
\begin{columns}
\column{.6\textwidth}
\centering
	\includegraphics[width = \textwidth, align = c]{Fig10-7}

\column{.35\textwidth}
\centering
\InstructorList{
  \item 192x179 image of tiger
	\item Convolution filters are 15x15 images with mostly zero (black) narrow strip of ones (white)
	\item Highlights areas of the tiger that resemble the filter
}

\end{columns}

\end{frame}
%--- Next Frame ---%
\begin{frame}[t]{More notes on convolution}

\begin{columns}
\column{.45\textwidth}
\centering
\InstructorList{
  \item Filters idea from image processing
	\item Novelty in DL is that the choice of filters are learned
	\item Filter weights: one hidden unit per pixel in convolved image; but more constrained than that
}

\column{.45\textwidth}
\centering
\InstructorList{
  \item Image is in color, so 3 channels
	\item This means the convolution filter will have three channels, one per color, each of dimension 3x3
	\item Results are summed to form a 2-dimensional output. Color is "forgotton" at this point
	\item $K$ different filters here means $K$ 2-d output maps, which can then be treated as channels
	\item Applying ReLU to convolved image. Sometimes viewed as separate layer in CNN, and then called a detector layer.
}

\end{columns}

\end{frame}
%--- Next Frame ---%

\begin{frame}[t]{Pooling layers}

\begin{columns}
\column{.45\textwidth}
\centering
\InstructorList{
  \item Condense a large image into a smaller summary image
	\item Lots of options, but max pooling is common:
	\item For each 2x2 non-overlapping block of pixels, get a single new pixel with the maximum value from the block
	\item Provides location invariance: so long as large value in one of the four pixels, the whole block registers as a large value
}

\column{.45\textwidth}
\centering
\includegraphics[width = \textwidth, align = c]{Eq10-3-2.png}
\end{columns}

\end{frame}
%--- Next Frame ---%

\begin{frame}[t]{Putting it together to make a CNN}

\includegraphics[width = \textwidth, align = c]{Fig10-8}
\begin{columns}
\column{.45\textwidth}
\centering
\footnotesize
\InstructorList{
  \item Start with three channels
	\item get a new channel at first hidden layer for each filter. Padding to keep the same size
	\item Next is max pool layer,
	\item convolve \& pool several more times
	\item Since size of image decreased, we increase number of filters in the next convolve layer to compensate
	% \item Can repeat several convolves before pool, this effectively increases the dimension fo the feature
}

\column{.45\textwidth}
\centering
\footnotesize
\InstructorList{
  \item After this, the maps are flattened (treated as separate units )
	\item Fed into 1 or more fully-connected layers
	\item Last is softmax activation for the 100 classes
}

\end{columns}

\end{frame}
%--- Next Frame ---%



%======================


\begin{frame}[t]{TL;DR}

\begin{columns}
\column{.45\textwidth}
\centering
\textbf{Feed Forward Neural Net}
\includegraphics[width = .5\textwidth, align = c]{Fig10-4}
\includegraphics[width = \textwidth, align = c]{Eq10-2}

\begin{itemize}
	\item Combines input data using learned weights
	\item Linear combo of those to get output
	\item Sometimes softmax to get probability of classification
\end{itemize}

\column{.45\textwidth}
\centering
\textbf{CNN}
\includegraphics[width = \textwidth, align = c]{Fig10-8}

\begin{itemize}
	\item Specialized NN
	\item Gets next layer via
	\begin{itemize}
		\item Convolution layer
		\item Pooling Layer
		\item Fully connected layerw
	\end{itemize}
\end{itemize}

\end{columns}


\end{frame}
%--- Next Frame ---%

\begin{frame}[t]{Next time}


\begin{columns}
\column{.45\textwidth}
\centering

\textbf{Next time:}

\begin{itemize}
	\item More Neural nets:
	\item Remainder of 10.3: CNNs
	\item 10.4: Document classification
	\item 10.5: Recurrent Neural Nets
\end{itemize}

\column{.45\textwidth}
\centering

\textbf{This week:}

	\begin{itemize}
	\item Mon 4/18: 10.1-3
	\item Weds 4/20: 10.4-5
	\item Fri 4/22: Lab: 10.9
	\begin{itemize}
		\item HW \#10 Due Friday
	\end{itemize}
	\item Week of 4/25: Ch 12 Clustering
	\item Final: May 4, 5:45 - 7:45 PM, Erickson 226

	\end{itemize}
	\end{columns}

\end{frame}
%--- Next Frame ---%


\end{document}
